\chapter{Lab 1}

\section{评分标准}
\textbf{60\%验收+40\%报告}

验收部分十个题,每个同学问两个总共占验收分数的30\%，70\% 实验部分由 50\% 仿真和 20\% 上板组成。

\section{预设问题}

\begin{enumerate}
    \item \textbf{Q：}实验背景中提到一个设计上的“妥协”：为了抵消 \texttt{ready-valid} 握手机制消耗的一个周期，将指令存储器取出的指令直接传入ID阶段，并将EX阶段计算出的访存地址直接传到数据存储器。请分析，为什么这种做法在更复杂的处理器设计中是不可取的？它可能会引入哪些问题？

\textbf{A：}限制时钟频率、增加设计复杂度、不利于异常处理。

\item \textbf{Q：}假设一条 \texttt{beq} 指令在EX阶段判断跳转成功，\texttt{flush} 机制应该具体如何操作？

\textbf{A：}清空（或置为无效）ID/EX寄存器、清空（或置为无效）IF/ID寄存器、更新PC。

\item \textbf{Q：}不在段间寄存器中存储控制信号，而是让每个阶段（EX, MEM, WB）都根据当前在该阶段的指令码 (\texttt{instruction[31:0]}) 重新进行译码并生成控制信号，会有什么问题？

\textbf{A：}硬件资源冗余与浪费（每个阶段都要译码）、影响性能、破坏流水线平衡、破坏设计的模块化和简洁性。

\item \textbf{Q：}与单周期CPU相比，五级流水线设计主要提高了CPU的哪一项性能指标？它是通过什么方式实现这一点的？

\textbf{A：}主要提高了CPU的吞吐率，它通过并行处理的方式来实现这一点。

\item \textbf{Q：}请全面地分析有哪些不同来源的数据最终需要在WB阶段被写回目标寄存器？

\textbf{A：}ALU的计算结果、从内存加载的数据、当前PC+4的值、U-type指令的立即数。

\item \textbf{Q：}为什么 PC 的值也需要逐级通过流水线寄存器（从IF/ID一直传递到MEM/WB）进行传递？

\textbf{A：}流水线中同时有多条处于不同执行阶段的指令，必须确保每个阶段使用的PC值是与当前阶段指令相对应的那个。如 \texttt{jal}, \texttt{auipc} 指令都需要 PC 的值。

\item \textbf{Q：}流水寄存器的引入虽然解决了指令并行执行的问题，但也带来了额外的开销，分析流水线寄存器对CPU性能的正面和负面影响。

\textbf{A：}
\begin{itemize}
    \item \textbf{正面影响：}分段后的流水线CPU可以运行在比单周期CPU高得多的时钟频率上。
    \item \textbf{负面影响：}对于单条指令而言，它的延迟变大了。在单周期CPU中，一条指令在一个时钟周期内完成；而在五级流水线中，一条指令需要经过五个时钟周期才能执行完毕。此外，流水线寄存器本身会带来额外的路径延迟和硬件电路开销。
\end{itemize}

\item \textbf{Q：}为什么要引入这个设计上的“妥协”：将指令存储器取出的指令直接传入ID阶段，并将EX阶段计算出的访存地址直接传到数据存储器？

\textbf{A：}为了抵消 \texttt{ready-valid} 握手机制消耗的一个周期。

\item \textbf{Q：}设想不引入设计上提前传递数据的妥协和 \texttt{flush} 机制的结果，并提出一种解决方案。

\textbf{A：}错拍/错误指令被提交导致错误结果，解决方案为 \texttt{stall/forwarding}。

\item \textbf{Q：}单周期CPU的 \texttt{RegFile} 需要两个读端口和一个写端口。在我们的五级流水线CPU中，寄存器堆是否需要更多的端口？

\textbf{A：}否，因为在任何一个时钟周期里，流水线中只有一条指令处于ID阶段，另一条指令处于WB阶段。
\end{enumerate}